\mypart{Множества, отношения и~отображения}

\section{Множества}

\begin{definition}
	\tdef{Множество}~"--- это набор различимых объектов произвольной природы.
\end{definition}

Множество можно задать перечислением элементов или заданием какого-нибудь свойства: если объект удовлетворяет этому свойству, то он является элементом заданного множества; если нет~"--- то нет.

Если множество задаётся перечислением, то его элементы перечисляются через запятую и~заключаются в~фигурные скобки:
\[
	A =
	\set*{\square, *, \lozenge, \mathghost, t}
	.
\]
Здесь~$A$~"--- это множество, состоящее из элементов $\square$, $*$, $\lozenge$, $\mathghost$ и~$t$.

Во множестве, заданном перечислением, порядок элементов не важен; однако каждый элемент входит ровно один раз. (<<Множество>>, в~которое один и~тот же элемент входит несколько раз и~все эти вхождения считаются различными, называется \tdef{мультимножеством}.)

Чтобы узнать, принадлежит ли элемент~$x$ множеству~$A$, заданному перечислением, нужно посмотреть, есть ли этот элемент в~списке элементов:
\[
	\bigl(
		x \in A
	\bigr)
	:=
	\left( 
		\mathtextbox{%
			элемент~$x$ есть в~списке элементов множества~$A$%
		}%
	\right)
	.
\]

Другой способ задать множество~"--- это задать свойство, которому должны удовлетворять все элементы множества и~только они.  Записывается это так:
\[
	A =
	\set{x}{P(x)}
\]
"--* и~читается <<$A$~"--- это множество таких элементов~$x$, что выполняется условие~$P(x)$>>. При этом~$P(x)$~"--- это некоторое высказывание: для каждого отдельного элемента~$x$ оно или истинно, или ложно.

Вертикальная черта выполняет роль связки <<такой, что>> или <<обладающий свойством>>. Иногда вместо вертикальной черты ставят двоеточие:
\[
	A =
	\set*{x : P(x)}
	.
\]

Итак, синтаксис имеет вид
\[
	\set{
		\mathtextbox[2.0in]{%
			элементы, из которых составлено множество%
		}%
	}{
		\mathtextbox[2.5in]{%
			свойство, которому должны удовлетворять элементы множества%
		}%
	}
	.
\]

Элементом множества может быть другое множество: возьмём, к~примеру, множество, состоящее из стола, стула и~списка учеников:
\[
	\set*{
		\text{стол},
		\text{стул},
		\set*{
			\text{ученик}_1,
			\dots,
			\text{ученик}_n
		}
	}
	.
\]

Если элемент~$x$ принадлежит множеству~$A$, то пишем~$x \in A$.

Если множество~$A$ задано некоторым свойством~$P$, то сказать $x \in A$~"--- то
же самое, что сказать, что~$x$ обладает свойством~$P$:
\[
	\bigl(
		x \in A
	\bigr) :=
	\left(
		\mathtextbox[1.29in]{%
			выполняется свойство~$P(x)$%
		}%
	\right)
	.
\]

Будем считать, что все наши множества будут жить в~одном большом множестве,
называемом \tdef{универс\'{у}мом Гротенд\'{и}ка}. В~нём, неформально говоря,
содержатся все множества, которые нужны в~данном рассуждении.

\begin{definition}
	\label{def:Universum}
	\tdef{Универсум Гротендика}~$\GU$~"--- это непустое множество, такое, что:
	\begin{enumerate}[1)]
		\item если $x \in y$ и~$y \in \GU$, то $x \in \GU$;
		\item если $x \in \GU$ и~$y \in \GU$, то $\set*{x, y} \in \GU$;
		\item если $x \in \GU$, то $2^x \in \GU$;
		\item если $\set{x_i \in \GU}{i \in I}$~"--- набор элементов из~$\GU$
			и~$I \in \GU$, то
					$
						\bigcup\limits_{i \in I} x_i \in \GU
					$.
	\end{enumerate}
\end{definition}

\begin{definition}
	\label{def:subset}
	Говорят, что множество~$A$ \tdef{является подмножеством} множества~$B$,
	и~пишут $A \subset B$, если из того, что $x \in A$, следует, что $x \in B$
	(любой элемент, лежащий в~$A$, обязательно лежит в~$B$):
	\[
		(A \subset B) :=
		\bigl(
			(x \in A)
			\implies
			(x \in B)
		\bigr)
		.
	\]
\end{definition}

Подмножество может совпадать с~объемлющим множеством. Подмножество~$A \subset B$, не совпадающее с~$B$, называется \tdef{собственным}, а~совпадающее с~$B$~"--- \tdef{несобственным}.

\begin{exercise}
	Какие из следующих утверждений верны? Указание: используйте
	определения операций~$\in$ и~$\subset$ (также может понадобиться
	вспомнить определение импликации~$\Longrightarrow$).
	\begin{multicols}{4}
		\begin{enumerate}[1)]
			\item $\emptyset \in     \emptyset$;
			\item $\emptyset \subset \emptyset$;
			\item $\emptyset \in     \set*{\emptyset}$;
			\item $\emptyset \subset \set*{\emptyset}$;
			\item $\emptyset \in     A$;
			\item $\emptyset \subset A$;
			\item $a \in \set*{a, b, c}$;
			\item $a \in \set*{ \{ a, b \}, c}$.
		\end{enumerate}
	\end{multicols}
\end{exercise}

\subsection{Примеры множеств}

\begin{example}
	Один из важных примеров~"--- это \tdef{пустое множество}, которое
	обозначается символом $\emptyset$. Формально его можно определить так:
	\[
		\emptyset := \set{x}{x \neq x}
		.
	\]
	Пустое множество не содержит в~себе элементов, так как элементов, не равных
	самим себе, нет (для того, чтобы это строго доказать, надо ввести понятие
	равенства, но пока что будем пользоваться интуитивным представлением о~нём).
\end{example}

\begin{example}
	Ещё один пример~"--- натуральные числа~$\NN$. Строгое его определение приведём
	позже, а~пока что будем довольствоваться таким:
	\[
		\NN := \set*{1, 2, 3, \dots}
		.
	\]
	Можно определить (хотя наивно и~не вполне корректно)~$\NN$ как множество
	чисел, которые <<используются при счёте>>.  Иногда (например, во Франции)
	в~$\NN$ включают ещё и~$0$ (и~в~этом есть смысл, если определять натуральные
	числа, не апеллируя к~физическому миру%
	\footnote{%
		Натуральные числа можно определить как классы мощностей конечных множеств.
		Введём на множествах функцию~$S$: для множества~$A$ положим
		\[
			S(A) := A \cup \set*{A}.
		\]
		Далее положим по определению~%
		$0 := \abs{\emptyset}$,
		$1 := \abs{S(\emptyset)} = \abs{\set*{\emptyset}}$,
		$2 := \abs{S \bigl( S(\emptyset) \bigr)} =
			\abs{\set*{\emptyset, \set*{\emptyset}}}$,
		и~так далее. Фраза <<и~так далее>> формализуется с~помощью леммы Цорна; об
		этом можно почитать, к~примеру, в~курсе А.\,Л.\,Городенцева <<Алгебра~"---
		1>>.

		Также можно определить натуральные числа как классы эквивалентности конечных множеств:
		\begin{align*}
			0 &:= \emptyset;
			\\
			1 &:= \set*{0} = \set*{\emptyset};
			\\
			2 &:= \set*{0, 1} = \set*{\emptyset, \set*{\emptyset}},
			\\
			3 &:= \set*{0, 1, 2} =
			\set*{\emptyset, \set*{\emptyset}, \set*{\emptyset, \set*{\emptyset}}},
		\end{align*}
		и~так далее. Вообще, $n + 1 = n \cup \set*{n} = \set*{0, 1, \dots, n}$. Это
		построение \tdef{конечных ординалов}.
	}%
	). Будем обозначать натуральные числа, в~которых присутствует~$0$,
	символом~$\NN_0$:
	\[
		\NN_0 := \set*{0, 1, 2, 3, \dots}
		.
	\]
\end{example}

\begin{example}
	Целые числа~"--- это натуральные числа, ноль и~обратные к~ним по сложению:
	\[
		\ZZ := \set*{\dots, -3, -2, -1, 0, 1, 2, 3, \dots} =
		\set*{0, \pm 1, \pm 2, \pm 3, \dots}.
	\]
\end{example}

\begin{example}
	Рациональные числа~$\QQ$~"--- это дроби, в~числителе которых стоят целые
	числа, а~в~знаменателе~"--- натуральные, исключая нуль. При этом одинаковыми
	считаются дроби $a/b$ и~$c/d$, если $ad = bc$:
	\[
		\QQ :=
		\left.
			\set{\frac{a}{b}}{%
				\begin{aligned}
					a &\in \ZZ, \\
					b &\in \NN \setminus \set*{0}
				\end{aligned}
			}
		\middle/
			\sim
		\right.
		.
	\]
	Для того чтобы корректно определить рациональные числа и~понять запись выше,
	нам пока не хватает техники и~понятия фактормножества.
\end{example}

\begin{example}
	Вещественные, или действительные, числа~$\RR$~"--- это рациональные и
	иррациональные числа. Строгое определение вещественных чисел достаточно
	сложно (существует несколько равносильных способов их определения), но одно
	из главных их свойств~"--- это \emph{полнота}. Формулируя упрощённо, можно
	сказать, что полнота заключается в~том, что между любыми двумя точками
	из~$\RR$ есть бесконечно много точек. При этом точек в~любом непустом
	интервале~"--- к~примеру, в~$(0, 1)$~"--- <<столько же>>, сколько во
	всём~$\RR$.
\end{example}

\begin{example}
	Компл\'{е}ксные числа~$\CC$~"--- это числа вида~$a + ib$, где $a$ и~$b$~"---
	вещественные, а~$i$~"--- \tdef{мнимая единица} со свойством $i^2 = -1$.
\end{example}

\begin{example}
	Читать выражение <<$a \in A$>> можно по-разному.
	Например,
	\begin{itemize}
		\item выражение $x \in \RR$ можно прочитать как <<$x$ принадлежить
			множеству действительных чисел>>, а~можно~"--- как <<$x$~"---
			действительное число>>; 
		\item выражение $x \in \set{a \in \NN}{a \mathbin{\smash{\vdots}} 2}$ можно
			прочитать как <<$x$ принадлежит множеству натуральных чисел, таких, что
			каждое из них делится на~$2$>>, или как <<$x$~"--- натуральное число,
			которое делится на~$2$>>, или как <<$x$~"--- натуральное чётное число>>.
	\end{itemize}
\end{example}

\subsection{Операции над множествами}

Несколько отклоняясь от основного повествования, заметим, что во избежание парадоксов <<наивной>> теории множеств типа парадокса Рассела (<<Существует ли множество всех множеств?>>) будем считать, что все множества, с~которыми мы будем работать, находятся в~одном большом множестве~"--- универсуме~$U$.

Элементы множеств будем называть \tdef{точками}.

С~множествами можно производить операции, среди которых нам понадобятся объединение, пересечение, разность и~декартово произведение.

Пусть даны множества~$A$ и~$B$.
\begin{definition}
	\begin{itemize}
		\item \tdef{Пересечение} множеств~$A$ и~$B$~"--- это множество~$A \cap B$, состоящее из элементов, принадлежащих одновременно и~$A$, и~$B$:
			\begin{equation}
				\label{eq:cap}
				A \cap B :=
				\set{x \in U}{
					(x \in A) \wedge (x \in B)
				}
				.
			\end{equation}

		\item
			\tdef{Объединение} множеств~$A$ и~$B$~"--- это множество~$A \cup B$,
			состоящее из элементов, принадлежащих хотя бы одному из множеств~$A$
			или~$B$~"--- то есть или~$A$, или~$B$, или обоим сразу:
			\[
				A \cup B :=
				\set{x \in U}{
					(x \in A) \vee (x \in B)
				}
				.
			\]

		\item \tdef{Разностью} множеств~$A$ и~$B$ называется множество~$A \setminus B$, которое состоит из элементов, которые присутствуют в~$A$, но не присутствуют в~$B$:
			\begin{equation}
				\label{eq:minus}
				A \setminus B := 
				\set{x \in U}{
					(x \in A) \wedge (x \notin B)
				}
				.
			\end{equation}

		\item \tdef{Дополнением множества} $A$ называется множество, которое обозначается $A'$ или $\overline{A}$ и~состоит из тех элементов, которые не принадлежат~$A$:
			\begin{equation}
				\label{eq:completement}
				\overline{A} := \set{x \in U}{x \notin A}
				.
			\end{equation}
	\end{itemize}
\end{definition}

Если внимательно посмотреть на определения пересечения множеств~\eqref{eq:cap}, их разности~\eqref{eq:minus} и~дополнения~\eqref{eq:completement}, то можно увидеть, что разность множеств можно переписать так:
\[
	A \setminus B = A \cap \overline{B}.
\]

В~булевой алгебре есть правила де Моргана:
\begin{align*}
	\overline{a \wedge b} &= \overline{\mathstrut a} \vee   \overline{b};
	\\
	\overline{a \vee b}   &= \overline{\mathstrut a} \wedge \overline{b}.
\end{align*}

Для множеств существуют их аналоги:
\begin{align*}
	\overline{A \cap B} &= \overline{A} \cup \overline{B};
	\\
	\overline{A \cup B} &= \overline{A} \cap \overline{B}.
\end{align*}
Докажем, к~примеру, первое из них:
\begin{align*}
	\overline{A \cap B}
	&=
	\overline{
		\set{x \in U}{
			(x \in A) \wedge (x \in B)
		}
	}
	\nlinea{=}
	\set{x \in U}{
		\overline{
			(x \in A) \wedge (x \in B)
		}
	}
	\nlinea{=}
	\set{x \in U}{
		\overline{
			x \in A
		}
		\vee
		\overline{
			x \in B
		}
	}
	\nlinea{=}
	\set{x \in U}{
		\overline{
			x \in A
		}
	}
	\cup
	\set{x \in U}{
		\overline{
			x \in B
		}
	}
	\nlinea{=}
	\overline{
		\set{x \in U}{
				x \in A
		}
	}
	\cup
	\overline{
		\set{x \in U}{
			x \in B
		}
	}
	\nlinea{=}
	\overline{A} \cup \overline{B}
	.
\end{align*}

Далее будут операции, которые будут встречаться нам реже~"--- кроме, пожалуй, декартова произведения.
\begin{itemize}
	\item \tdef{Декартовым произведением} множеств~$A$ и~$B$ называется
		множество~$A \times B$, состоящее из упорядоченных пар элементов
		множеств~$A$ и~$B$, в~каждой из которых первым идёт элемент из~$A$, а
		вторым~"--- элемент из~$B$:
		\[
			A \times B :=
			\set{(a, b)}{(a \in A) \wedge (b \in B)}
			.
		\]

		Обрати внмание, что в~записи <<$(a, b)$>> круглые скобки означают, что нам
		важен порядок элементов: парочка $(a, b)$ не равна парочке~$(b, a)$.

		Точно так же можно определить декартово произведение множества на себя:
		\[
			A \times A :=
			\set{(a_1, a_2)}{(a_1 \in A) \wedge (a_2 \in A)}
			.
		\]
\end{itemize}

\begin{definition}
	% Множество $A$ называется \tdef{подмножеством} множества~$B$, если любая точка, лежащая в~$A$, лежит и~в~$B$:
	% \[
	% 	(x \in A) \implies (x \in B)
	% 	.
	% \]
	Множество~$A$ называется \tdef{подмножеством} множества $B$, если каждый
	элемент, лежащий в~$A$, лежит в~$B$, то есть если из истинности $x \in A$
	следует истинность $x \in B$.
\end{definition}

Если $A$~"--- подмножество $B$, то пишем $A \subseteq B$ или $A \subset B$, в
обоих случаях подразумевая, что случай $A = B$ не исключается. Если нужно явно
указать, что $A \neq B$, то пишем $A \subsetneq B$; в~этом случае $A$
называется \tdef{собственным подмножеством} множества~$B$.

Итак, два следующих высказывания эквивалентны:
\[
	\bigl(
		A \subset B
	\bigr)
	:=
	\bigl(
		(x \in A) \implies (x \in B)
	\bigr)
	.
\]

\subsection{Некоторые синтаксические нужности}

Пусть множества~$A$ и~$B$ заданы своими свойствами~$P$ и~$Q$ соответственно:
\[
	A := \set{x \in U}{P(x)}
	,
	\quad
	B := \set{x \in U}{Q(x)}
	.
\]

Тогда по определению
\begin{align*}
	\overline{A} &= \set[big]{x \in U}{\overline{P}(x)},
	\\
	A \cap B &=
	\set{x \in U}{P(x)} \cap
	\set{x \in U}{Q(x)} =
	\set{x \in U}{\bigl( P(x) \bigr) \wedge \bigl( Q(x) \bigr)},
	\\
	A \cup B &=
	\set{x \in U}{P(x)} \cup
	\set{x \in U}{Q(x)} =
	\set{x \in U}{\bigl( P(x) \bigr) \vee \bigl( Q(x) \bigr)},
	\\
	A \setminus B &=
	A \cap \overline{B} =
	\set[big]{x \in U}{P(x)} \cap
	\set[big]{x \in U}{\overline{Q}(x)} =
	\set[big]{x \in U}{\bigl( P(x) \bigr) \wedge \bigl( \overline{Q}(x) \bigr)}.
\end{align*}

Иногда нужно бывает определить пересечение или объединение нескольких множеств, индексированных числами или элементами другого множества. Если даны множества~$A_1$, \dots, $A_n$, то
\begin{align*}
 \bigcap_{i = 1}^{n} A_i
	&:=
	A_1 \cap \dots \cap A_n,
	\\
	\bigcup_{i = 1}^{n} A_i
	&:=
	A_1 \cup \dots \cup A_n.
\end{align*}
Вместо переменной~$i$, считающей индексы от~$1$ до~$n$, может быть любая другая, не участвующая в~остальном выражении.

Объединять и~пересекать можно не только конечные наборы множеств: если $I$~"--- \emph{произвольное} множество индексов (конечное или бесконечное), а~$\set{A_i}{i \in I}$~"--- набор множеств, индексированных множеством~$I$, то определим
\begin{align*}
	\bigcap_{i \in I} A_i &:=
	\set{x \in U}{x \in A_i \text{ при всех $i \in I$}};
	\\
	\bigcup_{i \in I} A_i &:=
	\set{x \in U}{x \in A_i \text{ хотя бы при одном $i \in I$}}.
\end{align*}
Между кванторами $\forall$ и~$\exists$ и~операциями $\cap$ и~$\cup$ существует прямая связь: с~использованием кванторов определения пересечения и~объединения множеств можно записать так:
\begin{align*}
	\bigcap_{i \in I} A_i &:=
	\set{x \in U}{\forall i \in I \holds x \in A_i};
	\\
	\bigcup_{i \in I} A_i &:=
	\set{x \in U}{\exists i \in I \suchthat x \in A_i}.
\end{align*}

Пусть $A$ --- множество, а~$P$ --- высказывание. Определим такие множества:
\begin{align*}
	\bigl(
		\forall x \in A
		\holds
		P(x)
	\bigr)
	&:=
	\bigl(
		\forall x \in \GU
		\holds
		(x \in A)
		\implies
		P(x)
	\bigr)
	,
	\refstepcounter{equation}
	\tag{\theequation.1}
	\label{eq:sets:forall:x_in_set}
	\\
	\bigl(
		\exists x \in A
		\holds
		P(x)
	\bigr)
	&:=
	\bigl(
		\exists x \in \GU
		\holds
		(x \in A)
		\wedge
		P(x)
	\bigr)
	.
	\tag{\theequation.2}
	\label{eq:sets:exists:x_in_set}
\end{align*}
Вместо $\forall x \in \GU$ и~$\exists x \in \GU$ можно было бы писать просто
$\forall x$ и~$\exists x$.

Посмотрим на то, как можно составить отрицания:
\begin{align*}
	\neg
	\bigl(
		\forall x \in A
		\holds
		P(x)
	\bigr)
	&:=
	\neg
	\bigl(
		\forall x
		\holds
		(x \in A)
		\implies
		P(x)
	\bigr)
	\nlinea{:=}
	\bigl(
		\exists x
		\holds
		\overline{
			(x \in A)
			\implies
			P(x)
		}
	\bigr)
	\nlinea{\invcoloneq}
	\bigl(
		\exists x
		\holds
		\overline{
			\overline{
				(x \in A)
			}
			\vee
			P(x)
		}
	\bigr)
	\nlinea{\invcoloneq}
	\bigl(
		\exists x
		\holds
		\overline{
			\overline{
				(x \in A)
			}
		}
		\wedge
		\overline{
			P(x)
		}
	\bigr)
	\nlinea{\invcoloneq}
	\bigl(
		\exists x
		\holds
			(x \in A)
		\wedge
		\overline{
			P(x)
		}
	\nlinea{:=}
	\bigl(
		\exists x \in A
		\holds
		\overline{
			P(x)
		}
	\bigr)
	,
	\\
	\neg
	\bigl(
		\exists x \in A
		\holds
		P(x)
	\bigr)
	&:=
	\bigl(
		\exists x
		\holds
		(x \in A)
		\wedge
		P(x)
	\bigr)
	\nlinea{:=}
	\bigl(
		\forall x
		\holds
		\overline{
			(x \in A)
			\wedge
			P(x)
		}
	\bigr)
	\nlinea{\invcoloneq}
	\bigl(
		\forall x
		\holds
		\overline{
			(x \in A)
		}
		\vee
		\overline{
			P(x)
		}
	\bigr)
	\nlinea{:=}
	\bigl(
		\forall x
		\holds
		(x \in A)
		\implies
		\overline{
			P(x)
		}
	\bigr)
	\nlinea{:=}
	\bigl(
		\forall x \in A
		\holds
		\overline{
			P(x)
		}
	\bigr)
	.
\end{align*}

\begin{exercise}
	Будет ли сохраняться это свойство, если поменять
	в~\eqref{eq:sets:forall:x_in_set} и~\eqref{eq:sets:exists:x_in_set}
	местами ${\implies}$ и~$\wedge$?
\end{exercise}


\section{Отношения и~отображения}
\begin{definition}
	\tdef{Отношением~$R$} на множествах~$A$ и~$B$ называется подмножество их
	декартова произведения:
	\[
		R \subseteq A \times B
		.
	\]
\end{definition}

Обрати внимание, что не каждое подмножество $R \subset A \times B$ можно
представить в~виде произведения подмножеств $U \subset A$ и~$W \subset B$.

\begin{definition}
	\tdef{Отображением}, или \tdef{функцией}, $\phi \colon A \to B$ называется
	отношение на~$A \times B$, удовлетворяющее двум условиям:
	\begin{description}
		\item[всюду определённость:] <<область определения>> $\phi$~"--- всё
			множество~$A$ целиком, то есть для любой точки~$a \in A$ должна
			существовать точка~$\phi(a) \in B$:
			\[
				\forall a \in A\ 
				\exists \phi(a) \in B
			\]
			(у~каждой точки есть образ);
		\item[функциональность:] отображение $\phi$ не должно <<расклеивать
			точки>>, то есть одной точке $a \in A$ должно соответствовать не более
			одной точки из~$B$. Это можно записать так:
			\[
				(a_1 = a_2) \implies
				\bigl( \phi(a_1) = \phi(a_2) \bigr)
				,
			\]
			или, чуть более осмысленно, развернув импликацию и~навесив на обе части
			отрицание,
			\[
				\bigl( \phi(a_1) \neq \phi(a_2) \bigr)
				\implies
				(a_1 \neq a_2)
			\]
			(прообразы двух различных точек не совпадают).
	\end{description}
\end{definition}

Иными словами, отображение $\phi \colon A \to B$~"--- это отношение, которое
переводит \emph{каждую} точку из~$A$ в~\emph{ровно одну} точку из~$B$:
\[
	\forall a \in A\ 
	\exists! \phi(a) \in B
	.
\]
При этом вполне может оказаться так, что две различные точки $a_1 \neq a_2$
из~$A$ отображение~$\phi$ отправит в~одну точку $\phi(a_1) = \phi(a_2)$
(<<склеивать точки>> не запрещено).

Если отображение $\phi \colon A \to B$ переводит точку $a$ в~точку~$\phi(a)$,
то пишут $a \mapsto \phi(a)$ (обрати внимание на вертикальную черту слева у
стрелочки). Для указания множеств, на которых действует отображение,
используется <<обычная>> стрелка~$\to$, а~для указания действия отображения на
\emph{элементах} множества используется <<стрелочка с~чертой>>~$\mapsto$. Такое
обозначение удобно для задания функции, особенно если не хочется придумывать
имя:
\begin{align*}
	x &\mapsto \sqrt[3]{x};
	&
	2y &\mapsto y;
	\\
	u^2 &\mapsto u^3;
	&
	w &\mapsto \ln{\cos{w}}.
\end{align*}
В~более привычных обозначениях здесь записаны функции $f(x) := \sqrt[3]{x}$,
$g(y) := y/2$, $h(u) := u^{3/2}$, $s(w) := \ln{\cos{w}}$.

Также, если задано отображение $\phi \colon A \to B$, пишут для множеств
$A \xrightarrow{\phi} B$, а~для элементов~"--- $a \overset{\phi}{\longmapsto} \phi(a)$.

Отображение $\phi \colon A \to B$ задаётся не только правилом, сопоставляющим
элементу множества~$A$ элемент множества~$B$, но и~самими множествами~$A$
и~$B$. Поэтому отображения
\[
	\left( 
		\begin{gathered}
			% \NN &\xrightarrow{f} \NN
			\NN \xrightarrow{f} \NN
			\\
			% x &\overset{f}\longmapsto x^2
			x \overset{f}\longmapsto x^2
		\end{gathered}
	\right)
	\quad
	\text{и}
	\quad
	\left( 
		\begin{gathered}
			\RR \xrightarrow{g} \RR
			\\
			x \overset{g}\longmapsto x^2
		\end{gathered}
	\right)
\]
разные, хотя и~задаются одной формулой.

Для отображения $\phi \colon A \to B$ и~подмножества~$W \subset A$ определим
\[
	\phi(W) := \set{\phi(w)}{w \in W} \subset B
\]
--- подмножество в~$B$, в~которое~$\phi$ переводит множество~$W$.

Важный частный случай~"--- это $\phi(A)$~"--- множество точек в~$B$, имеющих прообраз в~$A$.

Если задано отображение~$\phi \colon A \to B$, то для каждого элемента~$b \in
B$ определено \emph{подмножество} (!) в~$A$
\[
	\phi^{-1}(b) :=
	\set{a \in A}{\phi(a) = b}
	,
\]
которое называется \tdef{прообразом} точки~$b$ при отображении~$\phi$ и~состоит
из таких точек~$a \in A$, которые переходят при отображении~$\phi$ в
фиксированную точку~$b \in B$. Обрати внимание, что прообраз может быть пустым,
то есть не содержать ни одного элемента (это означает, что никакая точка из~$A$
не переходит в~$b \in B$).

У~отображения $\phi \colon A \to B$ могут быть ещё два свойства:
\begin{description}
	\item[инъективность:] это означает, что $\phi$ <<не расклеивает>> точки, то есть переводит разные точки из~$A$ в~разные точки из~$B$:
		\[
			\bigl(
				a_1 \neq a_2
			\bigr)
			\implies
			\bigl(
				\phi(a_1) \neq \phi(a_2)
			\bigr)
			;
		\]
	\item[сюръективность:] это означает, что $\phi$ переводит~$A$ <<целиком>> в~$B$, <<не оставляя в~$B$ пустых мест>>, то есть что образ~$A$~"--- всё множество~$B$:
		\[
			\phi(A) = B
			.
		\]
\end{description}

\begin{definition}
	Отображение, одновременно инъективное и~сюръективное, называется
	\tdef{биективным}, или \tdef{взаимно однозначным}. (Иногда ещё говорят
	<<взаимно однозначное соответствие>>.)
\end{definition}

Биективное отображение $\phi \colon A \to B$ переводит каждую точку из~$A$ в~ровно одну точку из~$B$, причём у~каждой точки из~$B$ есть прообраз, состоящий ровно из одной точки из~$A$.
